\begin{remark}
\label{editorial-source-remark-intersection-powers-ideal}
Lemma \ref{editorial-source-lemma-intersect-powers-ideal-module-zero} in particular
gives $\bigcap_{n\geq0}I^n=(0)$ for a proper ideal $I$ in a
Noetherian local ring. More generally let $R$ be Noetherian,
let $I\subset R$ be an ideal, and retain an original element
$f\in\bigcap_{n\geq0}I^n$. Apply the proof of
Lemma \ref{editorial-source-lemma-intersection-powers-ideal-module} with $M=R$.
It gives a single $g\in1+I$ annihilating the whole ideal
$\bigcap_{n\geq0}I^n$, so in particular $gf=0$.
Every prime containing $I$ avoids $g$, hence
$V(I)\subset D(g)$. In $R_g$, the equality $gf=0$ implies
$f/1=0$. Thus the original $f$, and indeed every member of the
intersection, is zero on the same principal open neighbourhood
$D(g)$ of $V(I)$, with the coefficient condition $g\equiv1\pmod I$.
When $I=R$, the closed set $V(I)$ is empty and $g=0$ is
allowed; no nonempty neighbourhood or nonzero denominator is
asserted in that case.
\end{remark}